% TOP_PROBLEM: {"id": 1503, "title": "Algebraic Number Normality", "category_id": 1, "category": {"id": 1, "name": "number_theory", "display_name": "Number Theory", "description": "Properties of integers, prime numbers, Diophantine equations.", "slug": "number-theory", "order_index": 1, "created_at": "2026-07-31T15:26:25.670Z"}, "status": "open"}
% FIRST_POSTED: null
% ENTRY_KIND: statement_only
% Imported from: batch08.tex; UnsolvedMath ID 1503
\documentclass[11pt]{article}
\usepackage[margin=25mm]{geometry}
\usepackage{fontspec}
\setmainfont{Latin Modern Roman}
\usepackage{amsmath,amssymb,mathtools}
\usepackage{unicode-math}
\setmathfont{Latin Modern Math}
\usepackage{xurl,hyperref}
\hypersetup{colorlinks=true,urlcolor=blue}
\setcounter{secnumdepth}{0}
\setlength{\parindent}{0pt}
\setlength{\parskip}{5pt}
\setlength{\emergencystretch}{3em}
\newcommand{\sref}[2]{\hyperlink{source:#1:#2}{[S#2]}}
\newcommand{\cataloguescope}[1]{\paragraph{Recorded target.}#1}
\begin{document}
% UnsolvedMath ID 1503; source catalogue code NUM-009
\setcounter{section}{1502}
\section{Algebraic Number Normality}\label{1503}
{\small\sffamily UnsolvedMath ID 1503 · Number Theory}
\subsection{Definitions and mathematical statement}

\paragraph{Shared notation.}$\mathbb N=\{1,2,\ldots\}$, and $\mathbb Z,\mathbb Q,\mathbb R,\mathbb C$ denote the integers, rationals, reals and complex numbers. Any other conventions are specified locally.


For a real number \(x\), let \(\{x\}=x-\lfloor x\rfloor\)
be its fractional part. For an integer base \(b\ge2\), write
\(\{x\}=\sum_{j\ge1}d_jb^{-j}\), with
\(d_j\in\{0,\ldots,b-1\}\), choosing the expansion that is not
eventually all \(b-1\). For a word
\(w=(w_1,\ldots,w_\ell)\in\{0,\ldots,b-1\}^\ell\), define
\[
 A_{b,N}(x,w)=\#\{1\le j\le N:
       (d_j,\ldots,d_{j+\ell-1})=w\}.
\]
The number \(x\) is normal in base \(b\) if
\(\lim_{N\to\infty}A_{b,N}(x,w)/N=b^{-\ell}\) for every
\(\ell\ge1\) and every word \(w\). It is normal, without a base
qualification, if it is normal in every integer base \(b\ge2\).
A real number is algebraic if it is a root of a nonzero polynomial
with integer coefficients.
\paragraph{Statement recorded in the supplied source.}
Is every real algebraic irrational number normal? Explicitly,
for every \(\alpha\in\mathbb R\setminus\mathbb Q\) algebraic over
\(\mathbb Q\), every integer \(b\ge2\), every \(\ell\ge1\), and every
\(w\in\{0,\ldots,b-1\}^\ell\), is
\[
                \lim_{N\to\infty}\frac{A_{b,N}(\alpha,w)}N=b^{-\ell}?
\] \sref{1503}{2}
\subsection{Short English statement}
Does the expansion of every real algebraic irrational have the expected frequency for every finite digit block, in every integer base?
\subsection{Sources}
\begin{description}
\item[\hypertarget{source:1503:1}{S1}] ULAM.AI and the UnsolvedMath Contributors, \emph{UnsolvedMath: A Curated Collection of Open Mathematics Problems}, record 1503 (NUM-009). CC BY 4.0. \url{https://huggingface.co/datasets/ulamai/UnsolvedMath}.
\item[\hypertarget{source:1503:2}{S2}] Davar Khoshnevisan, \emph{Normal Numbers are Normal}, Clay Mathematics Institute Annual Report 2006, pp. 15 and 27--28, Section 2, block frequencies and Borel's question for algebraic irrational numbers. \url{https://www.claymath.org/library/annual_report/ar2006/06report_normalnumbers.pdf}.
\end{description}
\label{end:1503}
\end{document}
